% \section{Overview}
% \label{subsec:1.1}

This lecture introduces the ENPM702 course, covering its structure,
grading policies, assignment progression, and the course coding
conventions. The development environment (Ubuntu, Visual Studio Code,
CMake, and Git) and the Linux shell are covered in dedicated reading
modules that you should work through before the semester begins.

\subsection*{Learning Objectives}

\begin{remarkbox}{}
   By the end of this lecture, you will be able to:
   \begin{compactitem}
       \item Understand the course structure, objectives, and grading policies.
       \item Identify the software tools used in this course (Ubuntu, VSCode, ROS 2, Git, Valgrind, Doxygen).
       \item Explain how shell and \CC exercises contribute to your participation grade.
       \item Apply the course coding conventions, such as using \mintcpp{'\n'} instead of \mintcpp{std::endl}
       \item Know where to find the setup, Linux Shell, and VSCode and CMake reading modules.
   \end{compactitem}
\end{remarkbox}